\originalchapter{稳恒磁场与磁化}
\originalsection{电流与矢量势}
稳恒电流满足 $\nabla\cdot\jj=0$, 磁场满足 $\nabla\times\bb=\mu_0\jj$、$\nabla\cdot\bb=0$. 局域无单极磁场可写 $\bb=\nabla\times\vct A$. 取库仑规范 $\nabla\cdot\vct A=0$, 得 $\nabla^2\vct A=-\mu_0\jj$, 进而
\[
\vct A(\rr)=\frac{\mu_0}{4\pi}\int\frac{\jj(\rr')}{|\rr-\rr'|}\dd^3x',\qquad
\bb(\rr)=\frac{\mu_0}{4\pi}\int\frac{\jj(\rr')\times(\rr-\rr')}{|\rr-\rr'|^3}\dd^3x'.
\]
第二式是第一式对观察坐标取旋度得到的毕奥--萨伐尔定律. 对细导线, $\jj\dd^3x'$ 换成 $I\dd\vct l'$. 安培环路定律总成立, 但只有足够对称时才能从单一环路积分确定局部场.

\begin{example}
半径 $R$ 的圆环载流 $I$, 从 $+z$ 方向看电流逆时针. 求轴线上距环心 $z$ 的磁场.
\end{example}
\begin{solution}
对称性使横向分量抵消, 轴向分量的分子是 $I R^2\dd\phi$. 于是 $B_z=\mu_0IR^2/[2(R^2+z^2)^{3/2}]$. 远处 $B_z\simeq\mu_0IR^2/(2|z|^3)$, 恰与磁偶极矩 $\mm=I\pi R^2\hat{\vct z}$ 的轴向远场一致.
\end{solution}

\originalsection{磁偶极与接触项}
对局域稳恒电流, $\int\jj\dd^3x=0$. 可以对 $\nabla'\cdot(x_i'\jj)=J_i$ 积分证明, 边界项由局域性消失. 势的首个非零远场项由磁偶极矩 $\mm=\frac12\int\rr'\times\jj\dd^3x'$ 决定,
\[
\vct A_\mathrm{dip}=\frac{\mu_0}{4\pi}\frac{\mm\times\rr}{r^3}.
\]
对其取旋度并用上一章的分布恒等式,
\[
\bb_\mathrm{dip}=\PV\frac{\mu_0}{4\pi r^3}[3(\mm\cdot\nn)\nn-\mm]
+\frac{2\mu_0}3\mm\delta^3(\rr).
\]
磁接触项的系数与电接触项不同, 源外的相似形状不能把两者在原点处等同. 小电流环在外场中受力 $\vct F=\nabla(\mm\cdot\bb)$、力矩 $\vct\tau=\mm\times\bb$, 这里 $\mm$ 固定且只微分外场. 磁能在固定外场、固定偶极矩模型下为 $-\mm\cdot\bb$.

\originalsection{物质中的磁场}
磁化强度 $\vct M$ 是单位体积磁偶极矩. 相应束缚体电流 $\jj_b=\nabla\times\vct M$, 束缚面电流 $\vct K_b=\vct M\times\nn$. 将这些从安培方程移走, 定义 $\vct H=\bb/\mu_0-\vct M$, 得 $\nabla\times\vct H=\jj_f$. 小药盒给 $B_{2n}=B_{1n}$, 小环路给 $\nn\times(\vct H_2-\vct H_1)=\vct K_f$.

均匀磁化球的面电流是 $\vct K=M\sin\theta\hat{\vct\phi}$. 球内 $\bb=2\mu_0\vct M/3$, 球外等价于 $\mm=(4\pi a^3/3)\vct M$ 的偶极场. 这一结果可通过球内均匀场、球外偶极场的法向连续和切向跳跃得到. $\vct H_\mathrm{in}=-\vct M/3$ 是退磁场; 它表明物体形状会改变内部磁响应.

线性各向同性模型 $\vct M=\chi_m\vct H$ 给 $\bb=\mu\vct H$, $\mu=\mu_0(1+\chi_m)$. 铁磁迟滞不能用一个常数 $\mu$ 完整描述. 本章依据 8.07 第10、11讲及作业7、8.
